\begin{receta}{Salmorejo}{20 min (+ frío)}
    
    \begin{minipage}[t]{0.35\textwidth}
        \section*{Ingredientes}
        \begin{itemize}[leftmargin=*]
            \item 1 kg de Tomates maduros
            \item 200g de Pan duro
            \item 150ml de Aceite de oliva
            \item 2 Dientes de ajo
            \item Vinagre
            \item Sal y un poco de agua
            \item \textbf{Guarnición:} Huevo cocido y jamón en tacos
        \end{itemize}
    \end{minipage}
    \hfill
    \begin{minipage}[t]{0.6\textwidth}
        \section*{Preparación}
        \begin{enumerate}
            \item \textbf{Preparar el pan:} Partir el pan duro con la mano en trozos (no hace falta que sean muy pequeños) y ponerlos en un bol grande.
            \item \textbf{Los tomates:} Cortar los tomates en trozos medianos (cada tomate dividido en unos 7 o 10 trozos).
            \item \textbf{Mezcla base:} Añadir al bol con el pan los tomates, un poco de agua, aceite, un buen chorreón de vinagre, sal y los ajos picados (retirando siempre el corazón para que no repita).
            \item \textbf{Triturar:} Pasar todo por la batidora hasta que quede una crema fina y homogénea.
            \item \textbf{Reposo:} Dejar enfriar en la nevera varias horas. Servir con huevo cocido picado y tacos de jamón por encima.
        \end{enumerate}
    \end{minipage}

    \vspace{1em}
    \begin{tcolorbox}[colback=orange!5, colframe=orange!50, title=Truco, size=small]
        Si quieres que el salmorejo quede más emulsionado y brillante, añade el aceite de oliva poco a poco al final mientras sigues batiendo.
    \end{tcolorbox}
\end{receta}